\chapter{Standard deviation is not the same as knowing the true value}
\section{A fixed mass and fluctuating measurements}
Suppose the completed drone has fixed mass $m$. A scale gives
\begin{equation}
y_k=m+b+\epsilon_k.
\end{equation}
$b$ is an unknown calibration offset and $\epsilon_k$ represents changing
measurement effects. Repeating readings estimates dispersion. It does not make
the aircraft's mass fluctuate or reveal the offset by itself.

For $N$ readings:
\begin{align}
\bar y&=\frac1N\sum_{k=1}^Ny_k,\\
s^2&=\frac1{N-1}\sum_{k=1}^N(y_k-\bar y)^2,\qquad s=\sqrt{s^2},\\
u_A(\bar y)&=\frac{s}{\sqrt N}\quad\text{if samples are independent}.
\end{align}
$s$ describes individual-reading dispersion; $u_A$ describes uncertainty of
their average from repeatability. Repeating a biased scale 10000 times reduces
$u_A$ but does not remove $b$. A characterized reference is needed to assess bias.

\section{Type A and Type B evaluation}
Type A means evaluated statistically from data. Type B means evaluated using
other information such as a calibration certificate, resolution or a justified
physical bound. These describe evaluation methods, not merely random versus
systematic physical effects.

For certificate expanded uncertainty $U$ and coverage factor $k$:
\begin{equation}
u_{\rm certificate}=U/k.
\end{equation}
For rounding interval width $q$, if error is reasonably uniform on $[-q/2,q/2]$:
\begin{align}
u_{\rm resolution}^2
&=\frac1q\int_{-q/2}^{q/2}x^2\dd x=\frac{q^2}{12},\\
u_{\rm resolution}&=q/\sqrt{12}.
\end{align}
A justified uniform bound $[-a,a]$ gives $u=a/\sqrt3$.
A triangular distribution on those bounds gives $u=a/\sqrt6$.
The distribution is part of the assumption, not something every tolerance label
automatically specifies.

\section{Combine uncertainty contributions}
For independent additive contributions:
\begin{equation}
u_c^2=u_A^2+u_{\rm certificate}^2+u_{\rm resolution}^2+\cdots.
\end{equation}
If repeated data already capture quantization variation, account for that overlap
rather than counting the same source twice. For correlated quantities keep their
covariance. In general, for $\vect y=f(\vect\theta)$:
\begin{align}
J_f&=\frac{\partial f}{\partial\vect\theta},\\
\Sigma_y&\approx J_f\Sigma_\theta J_f^T.
\end{align}
For one scalar:
\begin{equation}
u_y^2\approx\sum_i
\left(\frac{\partial f}{\partial\theta_i}\right)^2u_i^2
 +2\sum_{i<j}\frac{\partial f}{\partial\theta_i}
\frac{\partial f}{\partial\theta_j}\Cov(\theta_i,\theta_j).
\end{equation}
This is a first-order approximation. Strong nonlinearity may require propagation
of a justified joint distribution, not invented independent Gaussian inputs.

\subsection{Worked example: not aircraft measurements}
Invented teaching readings in grams are $1200,1201,1199,1200,1200$.
They give $\bar y=1200$ g, $s=.7071$ g and $u_A=.3162$ g.
If separate residual scale-calibration standard uncertainty is 1 g and a
separately justified resolution contribution is .2887 g:
\begin{equation}
u_c=\sqrt{.3162^2+1^2+.2887^2}\approx1.0878\unit g.
\end{equation}
An illustrative $k=2$ expanded uncertainty is 2.18 g. Its coverage interpretation
needs assumptions about distributions and degrees of freedom. This is not a measured
confidence interval for IntelQcopter.

\section{Correlated samples and changing conditions}
For stationary normalized lag correlations $\rho_j$, an approximate effective
sample count is
\begin{equation}
N_{\rm eff}\approx\frac{N}{1+2\sum_{j=1}^J(1-j/N)\rho_j}.
\end{equation}
Then $s/\sqrt{N_{\rm eff}}$ can be more appropriate. Lag window $J$ and
stationarity need judgment. Temperature drift, mounting stress and motion can
invalidate one stationary-noise calculation.

In practice the reference also has uncertainty. Estimate a correction and retain
its uncertainty:
\begin{equation}
y_{\rm corrected}=y_{\rm instrument}-\hat b.
\end{equation}
The corrected result's budget must include $u(\hat b)$. For a fitted gain, retain
coefficient covariance and residual behavior. Independent validation points check
predictions outside the points used to fit. A small fitting residual does not
establish infinitely precise truth.

\chapter{Practical identification of the aircraft}
\section{Records and measurement priority}
Blank CSV templates are under \texttt{measurements/templates}. They are empty,
not synthetic evidence. Record units, timestamps, instrument identity, calibration
status, sensor/ESC settings, mounting, temperature and battery voltage. Preserve
raw readings before filtering.

\begin{enumerate}
\item Completed takeoff mass and centre of mass.
\item Rotor locations and assembled inertia.
\item Static thrust/current/RPM versus command and battery voltage.
\item Command-to-thrust/RPM delay and transients near hover.
\item Sensor noise, bias, scale and temperature effects.
\item Electrical calibration, capacity, camera geometry and timing/communication.
\end{enumerate}
These procedures can be prepared without selecting a Pi/Jetson board now.

\section{Mass and centre of mass}
Weigh the full flight configuration, including ESCs, wiring, fasteners and payload.
Earlier neglect of small components does not justify excluding them from final
measurements. Check zero and an independently characterized reference mass near
the working range. Weigh, remove and replace the aircraft repeatedly to include
placement variation. Repeat after battery/payload changes.

For component COM positions $\vect r_i$:
\begin{equation}
m=\sum_i m_i,\qquad \vect r_C=\frac{\sum_i m_i\vect r_i}{m}.
\end{equation}
For a stationary two-support balance at known positions $x_1,x_2$:
\begin{align}
R_1+R_2&=mg,\\
R_1x_1+R_2x_2&=mgx_C,\\
x_C&=\frac{R_1x_1+R_2x_2}{R_1+R_2}.
\end{align}
Repeat in the other direction. Unknown support moments or contact constraints
invalidate this simple force/moment balance.

\section{Inertia from components}
For a box with dimensions $l_x,l_y,l_z$ about its own COM:
\begin{equation}
J_{\rm box}=\frac m{12}\diag(l_y^2+l_z^2,l_x^2+l_z^2,l_x^2+l_y^2).
\end{equation}
Rotate each tensor into common axes and shift it to the aircraft COM:
\begin{equation}
J=\sum_i\left[R_iJ_iR_i^T+
m_i\bigl((\vect d_i^T\vect d_i)I-\vect d_i\vect d_i^T\bigr)\right],
\quad \vect d_i=\vect r_i-\vect r_C.
\end{equation}
This is the matrix parallel-axis theorem. Approximate geometry still gives an
approximate tensor even when the algebra is exact.

\section{Inertia from a bifilar pendulum}
Two equal parallel strings have length $L$, half separation $a$.
For small torsion angle $\theta$ about the appropriate vertical COM axis,
geometry gives
\begin{equation}
\Delta h\approx\frac{a^2\theta^2}{2L}.
\end{equation}
Potential energy is $U=mg\Delta h$, hence
\begin{align}
\tau_{\rm restore}&=-\frac{\partial U}{\partial\theta}
=-\frac{mga^2}{L}\theta,\\
J\ddot\theta+\frac{mga^2}{L}\theta&=0,\\
\omega_n&=\sqrt{\frac{mga^2}{LJ}},\qquad
\mathcal T=\frac{2\pi}{\omega_n},\\
J&=\frac{mga^2\mathcal T^2}{4\pi^2L}.
\end{align}
Measure $\mathcal T=t_{\rm elapsed}/N_{\rm cycles}$ over several cycles; repeat
at several small amplitudes to check the approximation. Reorient the unpowered
assembly with a correctly aligned fixture for other axes. Coupled swinging is
not a valid torsion record. Remove props for these non-powered mechanical tests.

If a fixture rotates too, subtract its independently characterized inertia about
the same axis, not its period:
$J_{\rm aircraft}=J_{\rm total}-J_{\rm fixture}$.
For independent input uncertainties:
\begin{equation}
\left(\frac{u_J}{J}\right)^2\approx
\left(\frac{u_m}{m}\right)^2+
4\left(\frac{u_a}{a}\right)^2+
4\left(\frac{u_{\mathcal T}}{\mathcal T}\right)^2+
\left(\frac{u_L}{L}\right)^2.
\end{equation}
COM/axis alignment, fixture flexibility and non-small motion add model errors.

\section{Static thrust, current and RPM}
Use the actual motor, ESC and propeller type on a secured force fixture with a
guard/exclusion arrangement. Never hold a running drone/motor by hand. Check
airflow clearance empirically: increasing clearance should stop changing the
steady result appreciably.

\begin{enumerate}
\item Calibrate the entire force path using known forces; verify sign, lever
ratio, tare and preload.
\item Record zero command, terminal voltage and temperature.
\item At several commands, wait for a steady plateau; record force, DC current,
voltage and RPM where available.
\item Repeat increasing/decreasing commands to expose hysteresis, heating and sag.
\item Repeat near the expected hover thrust and at relevant pack voltages, within
verified component/fixture limits.
\end{enumerate}
A scale's displayed gram is not a newton:
$1$ gram mass-equivalent corresponds to $.00980665$ N.
Upward thrust can reduce or increase a reading depending on fixture geometry.
Use the calibrated signed change and mechanical ratio.

Fit $T(u,V)$ rather than imposing linearity. For a local quadratic:
\begin{align}
T(u)&=\beta_0+\beta_1u+\beta_2u^2,\\
T'(u_0)&=\beta_1+2\beta_2u_0,\\
4T(u_0,V_0)&=mg.
\end{align}
Use repeated points for scatter and independent points for validation.
With RPM measurements, $\Omega=2\pi\,\mathrm{RPM}/60$ permits a fit of
$T=k_T\Omega^2$; examine residuals and range of applicability.

\section{Response delay, time constant and model order}
A steady scale can have unknown internal filtering and insufficient bandwidth
for a 40 ms transient. Use a timestamped force/RPM chain with adequate mechanical,
analog and digital bandwidth. A 500--1000 Hz logger gives many samples across
40 ms, but does not itself prove a fast force sensor.

For an accepted first-order model with delay $t_d$:
\begin{equation}
T(t)=T_0+(T_\infty-T_0)\left(1-e^{-(t-t_d)/\tau_m}\right),\quad t\ge t_d.
\end{equation}
At $t_d+\tau_m$:
\begin{equation}
\frac{T-T_0}{T_\infty-T_0}=1-e^{-1}\approx.63212.
\end{equation}
Use this for an initial estimate, then fit the complete transient. Capture repeated
small up/down steps near hover at several voltages with actual command timestamps.
If residuals reveal oscillation, two timescales or asymmetric rise/fall, retain a
higher-order/nonlinear model. Separate ESC/transport delay from mechanical lag.

\section{Yaw torque}
A torque fixture can measure reaction force $F$ at a known tangential lever $\ell$:
\begin{align}
Q&=\ell F,\\
u_Q^2&\approx F^2u_\ell^2+\ell^2u_F^2+2\ell F\Cov(\ell,F).
\end{align}
Record simultaneous thrust/RPM and fit $Q(\Omega)$ or $Q(T)$.
The current $\kappa=.015$ m is assumed, not measured. Isolate reaction torque
from off-centre thrust, support friction and other fixture moments.

\section{IMU bias, noise, scale and temperature}
Log fresh timestamps, axes, physical units, full-scale/DLPF/sample settings and
temperature on a fixed support with motors unpowered. Observe warm-up trends;
do not invent a universal fixed warm-up time.

\begin{enumerate}
\item Stationary gyro mean estimates zero-rate offset, not gyro scale factor.
\item Residuals around that mean estimate sample noise under those settings.
\item Six accelerometer faces estimate diagonal gain/offset.
\item Known reference angle/rate motion is needed for gyro scale/cross-axis error.
\item Repeat at stable temperatures; fit temperature behavior only if data support it.
\end{enumerate}
For $b_i(T)=a_i+c_i(T-T_{\rm ref})$, retain coefficient covariance and residuals.
Do not assign a slope merely because bias is temperature dependent.

One timescale statistic is Allan deviation. For successive non-overlapping
cluster means of duration $\tau$:
\begin{equation}
\sigma_A^2(\tau)=\frac1{2(M-1)}
\sum_{j=1}^{M-1}(\bar y_{j+1}(\tau)-\bar y_j(\tau))^2.
\end{equation}
Its dependence on $\tau$ helps expose white noise and drift, but needs long enough
records and does not automatically identify every process coefficient.

\section{Height, range and camera measurements}
\textbf{Barometer:} log compensated pressure at fixed height and temperature,
then validate several independently measured height differences. Distinguish
slow reference/weather drift from white noise.

\textbf{Lidar:} use independently measured target distances, multiple surfaces,
angles and lighting; preserve invalid returns. Measure beam direction and
sensor-to-body lever arm. Slant distance is not vertical height without geometry.

\textbf{Camera:} record identity/resolution/focus and measured target dimensions.
Capture varied target poses, fit intrinsics/distortion, validate on independent
poses and measure camera-to-body extrinsics. Reprojection error is not
automatically metric obstacle-range error. Arbitrary unknown object scale is not
observable from one image alone.

\section{Voltage, current, capacity and timing}
\textbf{Voltage:} compare the divider/ADC to a characterized meter at several
voltages. Fit gain/offset and keep reference/fit uncertainty.
\textbf{Current:} compare with a suitable DC reference over the intended range;
measure zero, sign, gain and temperature effects. A voltage divider alone does
not provide current.

\textbf{Capacity:} use an appropriate controlled charger and discharger, and integrate
\begin{equation}
C_{\rm measured}=\int I\dd t/3600
\end{equation}
under stated temperature, current and cutoff conditions.
Do not measure usable capacity by flying until the motors stop. Pack labels
are not a complete capacity/available-power model.
SOC is an uncertain estimate: crossing 20 percent does not prove exactly 20
percent of physical usable capacity remains. Starting charge, effective capacity,
current calibration and temperature affect that estimate.

\textbf{Timing:} record acquisition, reception, computation and actuation
timestamps. A timing GPIO plus suitable instrument can measure local execution.
One-way distributed latency needs synchronized clocks or another method that
separates latency from clock offset. Keep distributions and worst observed values.

\section{How measurements improve gains}
Replace the normalized surrogate $T_{\max}$ by measured hover slope $g_T=T'(u_0)$:
\begin{equation}
b_\phi=\frac{4ag_T}{J_x},\qquad b_h=\frac{4g_T}{m}.
\end{equation}
For fixed desired bandwidth, rate gains scale as $J_x/(ag_T)$:
\begin{equation}
\frac{\delta K_p}{K_p}\approx
\frac{\delta J_x}{J_x}-\frac{\delta a}{a}-\frac{\delta g_T}{g_T}.
\end{equation}
Propagate their joint covariance, including shared measurement sources.
Sensor data inform observation covariance $R$ and bias/process assumptions for
$Q$. A mass standard deviation cannot be copied into a sensor matrix with
different physical units.
